\section{Correspondence with Empirical Alpha Routes}

The compact-fiber derivation gives a structural electromagnetic bridge
coupling
\[
\alpha_F=\frac1{x_{\rm fiber}},
\]
where \(x_{\rm fiber}\) is the positive root of the compact-fiber bridge
polynomial. The physical identification
\[
\alpha_F=\alpha_{\rm low-energy}
\]
rests on the role of this coupling in atomic binding: leading
material-radiative binding contains two bridge vertices and therefore
carries the scale
\[
\alpha_F^2.
\]

Empirical determinations of the fine-structure constant do not all
measure a single primitive object directly. They infer the low-energy
electromagnetic coupling through different physical systems: atomic
binding, spectroscopy, atom recoil, magnetic moments, and electrical
metrology
\cite{NIST_CODATA_2022_alpha_inv,Parker2018_alpha_Cs,
Morel2020_alpha_Rb,Aoyama2019_electron_gminus2_review,
BIPM_SI_Brochure_9th}. The compact-fiber result should therefore be
compared to empirical alpha science at the level of route structure. The
fixed-point theorem supplies the structural coupling \(\alpha_F\); the
empirical routes read that coupling through route-specific dimensional
anchors, correction schemes, and operational conventions. The general
compact-fiber readout and matching-correction layer is treated separately
in \cite{Wells2026FiniteReadout}.

\subsection{Atomic binding and spectroscopy}

The most direct correspondence is the atomic binding route. In
conventional atomic physics, the Hartree energy is
\[
E_h=\alpha^2m_ec^2,
\]
and the Rydberg constant is
\[
R_\infty=\frac{\alpha^2m_ec}{2h}.
\]
Both formulas place the squared electromagnetic coupling at the leading
Coulomb scale:
\[
\text{atomic binding scale}\propto \alpha^2.
\]

The compact-fiber source-response structure gives the corresponding
relation
\[
E_{\rm atomic,F}
=
C_{\rm mass/readout}\alpha_F^2,
\]
where \(C_{\rm mass/readout}\) contains dimensional mass, action, and
finite-readout matching factors \cite{Wells2026FiniteReadout}. Thus the
leading-order correspondence is
\[
\alpha^2
\quad
\leftrightarrow
\quad
\alpha_F^2.
\]

This makes the Hartree--Rydberg route the strongest empirical bridge for
the present result. The same squared-coupling role appears on both sides.
What remains external to the alpha fixed-point theorem is not the
electromagnetic coupling itself, but the dimensional anchoring and
correction stack: mass readout, action readout, reduced-mass corrections,
nuclear-size corrections, and the conversion from structural energy scale
to conventional frequency or length units.

Spectroscopic routes have the same leading structure. For hydrogenic
systems, the gross scale is controlled by expressions equivalent to
\[
E_n\sim -\frac{\mu c^2\alpha^2}{2n^2},
\]
with additional contributions from fine structure, Lamb shifts, hyperfine
structure, finite nuclear size, recoil, and system-specific corrections.
The compact-fiber correspondence is therefore strongest at the leading
Coulomb scale, where spectroscopy reads \(\alpha_F^2\) through the same
two-vertex material-radiative source-response structure. Higher-precision
spectroscopic extraction is a correction-route problem, not part of the
alpha fixed-point theorem itself.

\subsection{\texorpdfstring{Recoil and \(h/m\)-based determinations}{Recoil and h/m-based determinations}}

Atom recoil determinations infer \(\alpha\) through relations involving
recoil measurements, mass ratios, the Rydberg constant, and the action
anchor
\[
h.
\]
These routes include high-precision cesium and rubidium recoil
determinations and reviews of the \(h/m\) route to \(\alpha\)
\cite{Parker2018_alpha_Cs,Morel2020_alpha_Rb,
Clade2019_alpha_h_over_m_review}. Schematically, such routes combine
\[
R_\infty,
\qquad
h/m,
\qquad
m_e,
\qquad
c,
\]
together with measured mass ratios and recoil observables. They are
therefore not measurements of a bare coupling alone; they infer the
coupling through an action/mass/readout chain.

In compact-fiber terms, the route decomposes into two structural layers.
The alpha fixed-point theorem supplies
\[
\alpha_F,
\]
while the compact-fiber Planck/action and gravitational-coupling analysis
supplies separate structural content for the action anchor
\[
h
\]
\cite{Wells2026PlanckGravity}. In the notation of that work, the
Planck/action calculation uses the effective structural factor
\[
R_{\rm eff}
=
\frac{R_1}{(1+m+e)^{1/3}},
\]
together with the Larson/Nehru natural-unit anchor system.

Thus recoil and \(h/m\)-based routes now involve two distinct
compact-fiber theorem chains: the electromagnetic bridge coupling
\(\alpha_F\) from the present paper, and the Planck/action bridge from
the separate Planck/gravity analysis. This does not close the recoil
route. A complete compact-fiber treatment would still require structural
mass ratios, recoil kinematics, reduced-mass corrections, and
route-specific metrological conventions. The significance is narrower:
these routes are no longer merely single-input correspondences to
\(\alpha_F\); they provide a setting in which separate compact-fiber
structures for electromagnetic coupling and action enter the same
empirical extraction pipeline.

\subsection{Magnetic-moment and electrical-metrology routes}

The electron anomalous magnetic moment route extracts \(\alpha\) through
the QED radiative expansion for
\[
a_e=\frac{g-2}{2}.
\]
In standard form, this expansion is expressed in powers of
\[
\frac{\alpha}{\pi},
\]
with higher-order hadronic, electroweak, and mass-dependent
contributions entering at the precision frontier
\cite{Aoyama2019_electron_gminus2_review,Aoyama2012_tenth_order_QED}.

The compact-fiber correspondence is that \(g-2\) is an amplitude-level
radiative readout of the same low-energy electromagnetic coupling. The
present paper does not derive the QED vertex expansion, the coefficient
sequence, or the amplitude-level radiative normalization. A compact-fiber
derivation of this route would require an amplitude-level radiative
bridge formalism, a compact-fiber interpretation of vertex corrections,
and a derivation or reconstruction of the \(\alpha/\pi\) expansion
structure. In the present paper, the \(g-2\) route is treated as a
correspondence to the same coupling rather than as an independently
derived measurement route.

Electrical-standard routes involve relations such as
\[
K_J=\frac{2e}{h}
\]
for the Josephson constant and
\[
R_K=\frac{h}{e^2}
\]
for the von Klitzing constant. These routes connect charge, action,
impedance, and electrical metrology to the fine-structure constant
\cite{BIPM_SI_Brochure_9th}.

The alpha theorem supplies the electromagnetic bridge coupling
\(\alpha_F\), and the Planck/action work supplies structural action
content for \(h\) \cite{Wells2026PlanckGravity}. Electrical-standard
determinations also require a compact-fiber account of charge readout,
impedance, and metrological convention. They therefore provide a partial
correspondence at present: they involve the same electromagnetic
coupling, but a complete compact-fiber derivation of the route belongs to
a broader charge/action/electrical-metrology program.

\subsection{Structural coverage by empirical route}

The empirical-route correspondence can be summarized by structural
coverage rather than by route names alone. The table records where
compact-fiber structural content is presently available and which route
layers remain imported or open.

\begin{center}
{\footnotesize
\setlength{\tabcolsep}{2pt}
\begin{tabularx}{\textwidth}{
  >{\hsize=0.65\hsize\raggedright\arraybackslash\hspace{0pt}}X
  >{\hsize=0.95\hsize\raggedright\arraybackslash\hspace{0pt}}X
  >{\hsize=1.25\hsize\raggedright\arraybackslash\hspace{0pt}}X
  >{\hsize=1.55\hsize\raggedright\arraybackslash\hspace{0pt}}X
  >{\hsize=0.60\hsize\raggedright\arraybackslash\hspace{0pt}}X}
\toprule
Route & Quantities used & Compact-fiber structural content & Remaining imports or open layers & Coverage \\
\midrule
Hartree--Rydberg & \(\alpha^2\), \(m_e\), \(c\) & \(\alpha_F\) from the present bridge equation & electron-mass readout, dimensional/unit conventions & coupling \\
Spectroscopy & \(\alpha^2\) plus transition corrections & \(\alpha_F\) at the leading Coulomb scale & fine structure, Lamb shift, hyperfine, recoil, finite-size, and route-specific corrections & coupling \\
Recoil / \(h/m\) & \(\alpha\), \(R_\infty\), \(h/m_X\), mass ratios & \(\alpha_F\) from this paper; Planck/action content & mass ratios, recoil kinematics, reduced-mass and route-specific correction pipeline & coupling + action \\
Electron \(g{-}2\) & QED expansion in \(\alpha/\pi\) & \(\alpha_F\) as the low-energy coupling entering the amplitude expansion & QED vertex expansion, coefficient sequence, hadronic/electroweak terms, amplitude-level compact-fiber readout & coupling \\
Electrical standards & \(e\), \(h\), \(R_K\), \(K_J\), impedance relations & \(\alpha_F\); Planck/action content & charge readout, impedance normalization, electrical-metrology conventions & coupling + action \\
\bottomrule
\end{tabularx}
}
\end{center}

This table is not a list of completed route derivations. It records where
the compact-fiber structural coupling and related compact-fiber theorem
paths enter existing empirical alpha science.

\subsection{Multi-route consistency}

The compact-fiber bridge coupling \(\alpha_F\) is derived without using
any route-specific empirical extraction of \(\alpha\) as an algebraic
input. Hartree/Rydberg spectroscopy, atom-recoil routes, electron
\(g{-}2\), and electrical-standard routes access the low-energy
electromagnetic coupling through different auxiliary structures:
electron mass, action, recoil kinematics, radiative vertex coefficients,
charge, and impedance conventions.

This route diversity strengthens the role identification because the
routes place the electromagnetic coupling in different auxiliary
structures: atomic binding, recoil/action chains, radiative vertex
expansions, and electrical metrology. It should not, however, be read as
a completed metrological closure. The present paper does not derive every
route-specific correction pipeline, nor does it assign a combined
statistical likelihood over the empirical route network. The point is
more limited: the same compact-fiber bridge coupling occupies the
coupling slot that these different operational routes infer. Agreement
with the empirical low-energy value is therefore a consistency check of
the role identification
\[
\alpha_F=\alpha_{\rm low-energy},
\]
not merely a comparison to one favored extraction route.

The important point is that the compact-fiber value is not identified
with the fine-structure constant merely because \(x_{\rm fiber}\) is
numerically close to the measured inverse value. The identification rests
on the structural role of \(\alpha_F\) as the material/radiative bridge
coupling and on the fact that the leading atomic scale reads that bridge
through \(\alpha_F^2\). The major empirical routes then appear as
different operational readouts of the same low-energy coupling, each with
its own dimensional anchors and correction structure.